%!TEX root = ../template.tex
%%%%%%%%%%%%%%%%%%%%%%%%%%%%%%%%%%%%%%%%%%%%%%%%%%%%%%%%%%%%%%%%%%%%
%% chapter2.tex
%% NOVA thesis document file
%%
%% Chapter with the template manual
%%%%%%%%%%%%%%%%%%%%%%%%%%%%%%%%%%%%%%%%%%%%%%%%%%%%%%%%%%%%%%%%%%%%

\typeout{NT FILE chapter2.tex}%

\chapter{Related Work}
\label{cha:related_work}

In this chapter, the relevant concepts and related work on consistency, replication, conflict resolution in distributed systems, and verification methods are presented. It starts by introducing consistency and replication models, clarifying the trade-offs between strong and weak consistency. The chapter then discusses how divergent replicas are handled, covering both ad-hoc conflict resolution strategies and principled approaches based on Conflict-free Replicated Data Types (CRDTs)~\cite{preguica2018overview}. Then, the synchronization models for CRDTs and the challenges that arise when composing and nesting replicated data structures are discussed. Finally, it addresses the role of formal verification in ensuring correctness, reviewing different verification approaches and the tools commonly used in practice, along with their capabilities and limitations.

\glsresetall

\section{Consistency models and Replication}
\label{sec:consistency}

Consistency models are commonly used to characterize the behavior of a system when multiple write and read operations are executed concurrently, defining the conditions under which different results can be observed. By establishing these rules, the model determines the level of consistency that applications can rely on. However, stronger consistency guarantees often come at the cost of reduced availability provided by the system, as established by the CAP theorem~\cite{cap-theorem}.

The CAP theorem states that a distributed system can not simultaneously provide strong consistency and high availability in the presence of network partitions. Since network partitions are unavoidable in practice, systems must adopt different trade-offs depending on the perpouse of the system.

Under strong consistency assumptions, it is considered that once an operation completes, all replicas reflect a uniform state, ensuring that clients observe the most recent write regardless of the node accessed. Alternatively, approaches that prioritize availability relax these guarantees to maintain availability during failures.

\subsection{Weak Consistency}
\label{sub:weak}

Systems that adopt a weak consistency model prioritize high availability and low latency when responding to client requests. In this model, client operations are executed on a small subset of replicas, which execute the operation locally and reply to the client before the update is propagated to the remaining replicas. This propagation occurs asynchronously, avoiding the need for coordination with a majority of replicas, resulting in lower latency and higher availability. However, the relaxed coordination requirements among replicas lead to divergencies, where different replicas can return different values for the same read operation. While these systems may eventually converge, the temporary divergence between replicas increases the complexity of reasoning about system behavior, making application development and debugging more challenging~\cite{preguica2018overview}.

\begin{description}
    \item[Eventual Consistency] is a weak consistency model that prioritizes availability over consistency guarantees. Under this model, replicas are allowed to diverge temporarily, with the only guarantee being that, in the absence of further updates, all replicas will eventually converge to a common state. During this convergence period, clients may observe out-of-date or inconsistent values across replicas.

          To ensure convergence, eventual consistent systems rely on additional mechanisms for resolving conflicting updates.
    \item[Causal Consistency] is one of the strongest weak consistency models that can be achieved while maintaining availability. This model requires the system to maintain causal dependencies between operations and guarantees that all clients observe causally related operations in an order that respects these dependencies.

          Two operations are causally related when the execution of one can influence the execution of the other. For example, in a social networking application, if a user creates a post and subsequently deletes it, the delete operation is causally dependent on the create operation and must be observed after it. Operations that are not causally related are considered concurrent and may be observed in different orders by different replicas. With this enforced consistent ordering of causally related operations while allowing flexibility in the ordering of concurrent ones, causal consistency provides higher consistency guarantees than eventual consistency without sacrificing availability~\cite{borrego2025convergence}.
    \item[Causal+ Consistency] extends causal consistency by ensuring that replicas converge to the same state even in the presence of concurrent updates. Although causal consistency preserves the ordering of causally related operations, it allows concurrent operations to be applied in different orders, potentially leading to divergent replica states. Causal+ consistency overcomes this issue by enforcing a convergent conflict resolution strategy, in which concurrent operations are merged using associative and commutative functions, guaranteeing eventual convergence while preserving causal dependencies~\cite{shapiro2011comprehensive}.
\end{description}

\subsection{Strong Consistency}
\label{sub:strong}

Systems that adopt a strong consistency model prioritize providing a single, well-defined view of the system state to all clients. In these systems, client operations are executed in a coordinated manner across replicas, ensuring that read operations always observe the effects of previously completed writes according to a well-specified ordering.

To enforce this guarantee, strong consistency models require tighter coordination among replicas, often involving synchronous communication and agreement protocols before operations are considered complete. While this coordination simplifies reasoning about system behavior and enables applications to rely on intuitive semantics, it typically comes at the cost of increased latency and reduced availability when compared to weakly consistent systems.

\begin{description}
    \item[Linearizability] is a strong consistency propertie that enforces a strict ordering of operations in replicated systems. It requires that each operation on a shared object appears to take effect at a single, indivisible point in time between its invocation and completion. Consequently, the outcome of all operations must be consistent with a global order that reflects the real-time order in which operations are issued by clients.

          This propertie guarantees that once a write operation completes, its effects are immediately visible to all subsequent operations, independently of the replica they access.
    \item[Serializability] is a consistency propertie that considers the execution of transactions, each composed of one or more operations that may access multiple objects. It guarantees that the execution of a set of concurrent transactions is equivalent to a serial execution, where transactions are totally ordered.

          With this propertie, all replicas observe transactions in the same order. However, this order is not required to respect real-time constraints. Therefore, the serialization order may differ from the temporal order in which transactions are issued according to a global wall clock.
    \item[Sequential Consistency] is a consistency propertie in which all clients observe operations as if they were executed in a single global order. This order is shared by all replicas, ensuring that every client perceives the same sequence of operations, even if this sequence does not correspond to the real-time order in which operations are issued.

          While the global execution order does not need to respect real-time constraints, it must preserve the program order of each individual client or process. That is, if a process issues an operation op1 before op2, all replicas must observe op1 preceding op2 in the global sequence.

          Sequential consistency operates with individual operations. For this reason, systems that guarantee serializability also satisfy sequential consistency, whereas the opposite is not true. By enforcing a common operation order across replicas, sequential consistency prevents clients from observing divergent states when accessing different replicas.
\end{description}

\subsection{Strong Eventual Consistency}
\label{sub:sec}

Strong Eventual Consistency characterizes a class of replication models in which replicas are allowed to execute updates independently and propagate them asynchronously, without relying on coordination or global ordering. The central guarantee of this model is that any replicas that have applied the same set of updates will deterministically converge to an identical state, regardless of the order in which those updates are delivered or executed~\cite{shapiro2011comprehensive}.

This model is typically defined by three properties. \texttt{Eventual Delivery}, ensures that all updates are eventually propagated to all correct replicas. \texttt{Strong Convergence} guarantees that replicas that have applied the same set of updates reach an identical state. \texttt{Causal Consistency} preserves the ordering of causally related updates. Together, these properties prevent permanent divergence while avoiding the coordination overhead required by strongly consistent models.

Strong Eventual Consistency directly addresses the limitations of weakly consistent systems by eliminating non-deterministic conflict resolution, while retaining high availability and low latency under network partitions. In practice, strong eventual consistency is achieved through Conflict-free Replicated Data Types, whose design ensures that concurrent updates can be merged deterministically, enabling convergence without coordination~\cite{preguica2018overview}.

\section{Conflict Resolution Strategies}
\label{sec:conflict}

In practice, many replicated systems adopt weak consistency models in order to achieve high availability and low latency, allowing replicas to diverge temporarily as updates are executed independently and propagated asynchronously. While this design choice improves system responsiveness, it introduces the possibility of conflicting updates, making conflict resolution a fundamental requirement to obtain a consistent state.

Conflict resolution strategies define how concurrent or divergent updates are reconciled once replicas exchange their states or operations. Historically, these strategies have been broadly classified into ad-hoc approaches, which resolve conflicts using application-specific or heuristic rules, and principled approaches based in formal properties, such as those employed by Conflict-free Replicated Data Types. These approaches have a direct impact on system semantics, predictability, and scalability.

Consequently, conflict resolution is closely linked to the underlying replication and consistency model, influencing not only correctness guarantees but also the degree of coordination required and the overall system behavior under concurrency and failures.

\subsection{Ad-hoc Strategies}
\label{sub:strategies}

Ad-hoc conflict resolution strategies resolve conflicts using simple, often application-specific rules, without providing formal guarantees of deterministic convergence in the presence of concurrent updates, as they may lead to undesirable behaviors.

\begin{description}
    \item[Last-Writer-Wins] is one of the most common conflict resolution policies. Under this strategy, concurrent updates are resolved by selecting the update with the most recent timestamp or highest value according to a predefined global ordering.

          Although Last-Writer-Wins is simple to implement and introduces minimal coordination overhead, it is prone to lost updates, as concurrent modifications may overwrite each other without being merged. For example, in a replicated shopping cart, concurrent additions of different items by multiple users may result in only one of the updates being reflected in the final state, leading to unintended data loss~\cite{shapiro2011comprehensive}.
    \item[Programmatic conflict resolution] delegates the responsibility of resolving conflicts to the application logic. This approach was popularized by systems such as Amazon Dynamo~\cite{dynamo}, where replicas may return multiple conflicting versions of the same object to the client, which is then responsible for reconciling them and writing back a resolved version. While this strategy provides flexibility and allows applications to define domain-specific resolution semantics, it places a significant load on developers. Designing correct and efficient merge logic is often complex, particularly in the presence of high concurrency, and errors in conflict resolution may lead to inconsistent or unintended application states~\cite{kohler2019localfirst}.
\end{description}

\subsection{Conflict-free Replicated Data Types}
\label{sub:crdts}

Conflict-free Replicated Data Types emerge as a rigorous alternative for addressing replica divergence without the drawbacks associated with ad-hoc conflict resolution strategies. Rather than resolving conflicts through heuristic or application-specific rules, CRDTs embed conflict resolution directly into the data type design, ensuring predictable and well-defined behavior under concurrent updates.

A CRDT is an abstract data type designed to be replicated across multiple nodes, allowing replicas to perform local updates without requiring coordination. Convergence is ensured through mathematical principles, such as the commutativity and associativity of operations or the use of state-based merge functions defined over join-semilattices~\cite{shapiro2011comprehensive}. Therefore, any two replicas that have applied the same set of updates are guaranteed to reach an identical state deterministically, regardless of message ordering or network delays.

CRDTs include a variety of classes designed for different application needs, including counters, registers, sets, and graphs. Some of these data types exhibit recursive structures, where operations may depend on nested or hierarchical elements, enabling the design of more complex CRDTs that maintain convergence guarantees even in recursive scenarios~\cite{baquero2017nested}.

In contrast to ad-hoc strategies, CRDTs ensure that all concurrent updates are reflected in the final state, avoiding unintended data loss. Since CRDTs do not rely on consensus protocols in the critical path of update execution, they exhibit high scalability and fault tolerance. Updates can be applied with low latency and propagated asynchronously, enabling systems based on CRDTs to remain available and responsive even in the presence of network partitions.

\subsection{Hybrid and Tunable Policies}
\label{sub:hybrid}

While standard CRDTs provide strong convergence guarantees, they typically enforce a fixed conflict resolution policy, such as add-wins or remove-wins semantics, which may not be suitable for all application scenarios. In practice, applications often require policies that depend on the context, for example, temporary disconnections may favor add-wins semantics to preserve availability, while explicit actions such as removals or bans may require remove-wins semantics to ensure they take effect. To address this limitation, hybrid and tunable approaches have been proposed~\cite{rijo2018building}, allowing conflict resolution semantics to be adapted for individual updates. Tunable CRDTs enable developers to select the desired resolution policy for each update, rather than committing to a single global policy.

\section{Synchronization Models of CRDTs}
\label{sec:sync_models}

The synchronization model defines the assumptions and guarantees required by a system to ensure the correct behavior of CRDTs in a replicated environment. In particular, it specifies how updates are propagated among replicas and under which conditions convergence is achieved, ensuring that all replicas that receive the same set of updates eventually reach an identical state. To this end, CRDTs rely on two fundamental synchronization models, state-based replication, in which replicas periodically exchange their full state, and operation-based replication, where individual updates are disseminated and applied at remote replicas~\cite{shapiro2011comprehensive}.

\subsection{Convergent Replicated Data Types}
\label{sub:cvrdts}

Convergent Replicated Data Types adopt a state-based synchronization model, in which replicas periodically propagate their entire local state to other replicas. Upon receiving a remote state, a replica integrates it with its own state through a deterministic merge function, ensuring that both replicas advance towards a common state. This approach decouples local updates from coordination, allowing replicas to apply operations independently and reconcile divergences through state exchange.

Convergence in CvRDTs is achieved by constraining how replica states evolve and how divergent states are reconciled. The set of possible states is partially ordered, allowing replicas to compare states in terms of the information they contain. Local updates are required to be monotonic, meaning that applying an update to a state always produces a new state that is greater than or equal to the previous one. When replicas exchange state, reconciliation is performed by a deterministic merge function that computes the least upper bound of the two states, preserving all information observed by either replica. This merge operation is necessarily commutative, associative, and idempotent, ensuring deterministic convergence despite message reordering, duplication, or frequency of state exchanges.

From a systems perspective, CvRDTs impose minimal requirements on the underlying network. They tolerate message loss, duplication, and reordering, assuming only that communication is eventually reliable and that replicas can re-establish connectivity after partitions. This robustness makes CvRDTs particularly well-suited for highly unreliable or environments subject to network partitions. However, this model may incur significant communication overhead, as entire states must be transmitted during synchronization, which can become inefficient for frequently updated objects.

\subsection{Commutative Replicated Data Types}
\label{sub:cmrdts}

Commutative Replicated Data Types adopt an operation-based synchronization model, in which replicas propagate individual update operations rather than exchanging their full state. Each update is executed locally and disseminated to other replicas, allowing changes to be applied incrementally and avoiding the transmission of large state payloads.

When an update is issued, it is first processed at the origin replica, where any necessary context is used to produce an operation that can be safely applied at other replicas. This operation is then disseminated to the remaining replicas, where it is applied to the local state of each one. Therefore, replicas can generate updates independently, while ensuring that all replicas eventually apply the same operation and reach an identical state.

Convergence in CmRDTs relies on the commutativity of update operations. Since replicas may receive and apply operations in different orders, all operations that can be executed concurrently are required to commute, ensuring that the order in which updates are applied does not affect the final state. Once all replicas have received the same set of operations, they deterministically converge to an identical state, without requiring global coordination or synchronization.

Compared to state-based CRDTs, CmRDTs place stronger requirements on the underlying communication channel. Operations must be reliably delivered to all replicas, typically with causal ordering guarantees, to ensure that concurrent updates commute as expected and replicas do not diverge. While this increases the complexity of the communication layer, it significantly reduces synchronization overhead and makes CmRDTs more efficient for objects with large state or high update rates.

\subsection{Composition and Nesting}
\label{sub:nesting}

While the properties of basic CRDTs are well understood, practical systems frequently require combining these primitives into more complex, hierarchical structures, such as maps containing nested objects, JSON documents, or directory trees. In these settings, simply nesting CRDTs is not sufficient, as the convergence guarantees of individual data types do not automatically extend to the composed structure. A challenge arises when concurrent operations affect different levels of the hierarchy, for instance, when an update modifies a nested element while a concurrent operation removes or replaces its parent.

These challenges have been studied in the context of replicated JSON documents~\cite{kleppmann2017json}, which are often regarded as a reference model for conflict resolution in deeply nested data. CRDT based approaches to JSON highlight that simple resolution policies, such as Last-Writer-Wins, are insufficient to preserve user intent in hierarchical structures, as they may discard meaningful concurrent updates. Consequently, specialized mechanisms are required to reconcile edits across multiple levels of the hierarchy while maintaining deterministic convergence.

To address these issues, several mechanisms for combining CRDTs have been proposed. One commonly adopted approach is recursive reset semantics, where removing an element from a collection implicitly resets or discards all nested state. More recent approaches generalize this idea by explicitly capturing the causal relationships between parent and child objects. In particular, nested operation-based CRDT models~\cite{baquero2017nested} introduce mechanisms that allow removals or resets to be applied selectively and recursively, based on causal information, ensuring that the semantics of the enclosing structure are preserved, providing deterministic convergence for complex, hierarchical CRDTs without introducing excessive metadata overhead.

\section{Formal Verification}
\label{sec:verification}

Formal verification of replicated data types is essential because it is difficult to reason about the many possible execution orders of concurrent operations in distributed systems. Subtle combinations of concurrent updates, message reordering, or partial failures can lead to unexpected behaviors that are hard to detect through testing alone. Therefore, informal reasoning or ad-hoc validation techniques are often insufficient to guarantee correctness~\cite{borrego2025convergence}.

Research in verification of replicated data types has evolved from manual, pen-and-paper proofs toward automated and deductive verification techniques. Early approaches relied heavily on human reasoning to establish convergence and correctness properties, while more recent work has focused on developing formal frameworks and tool-supported methods that provide stronger guarantees and better scalability~\cite{porre2023masses}.

\subsection{Manual Verification}
\label{sub:manual}

Early efforts to validate CRDTs relied primarily on manual, pen-and-paper proofs~\cite{shapiro2011comprehensive}. These proofs aimed to establish key properties such as convergence and correctness by reasoning about the abstract behavior of the data type under concurrency. From a theoretical perspective, this approach offers a high level of rigor, as it allows precise arguments to be constructed over well-defined formal models.

However, manual verification presents significant limitations. Constructing such proofs is time consuming and requires substantial expertise in formal methods, making the approach difficult to scale beyond a small number of data types. Moreover, these proofs often reason about an abstract specification of the algorithm, which may differ from its concrete implementation. Consequently, even when the specification is formally correct, errors may still occur in the implementation.

\subsection{Automated Verification}
\label{sub:automated}

To make formal verification more accessible to non-specialists, several automated verification approaches have been proposed, often based on high-level specification languages. One such example is VeriFx~\cite{porre2023masses}, a framework that allows replicated data types to be implemented and automatically verified for Strong Eventual Consistency using SMT solvers such as Z3. By abstracting away low-level proof details, these tools significantly reduce the complexity of verifying replicated data structures and have been successfully applied to a wide range of basic CRDTs.

Despite these advances, the literature identifies important limitations in existing automated verification tools. In particular, handling recursive algorithms remains challenging, as demonstrated by the inability of VeriFx to verify data types such as the Replicated Growable Array, whose correctness depends on recursive insertion logic. Furthermore, the verification of more complex objects, especially those relying on version vectors, can be unstable in practice. For instance, small and semantically irrelevant changes to a specification, such as renaming variables, may cause the solver to fail or abort the proof process, highlighting limitations in robustness and predictability.

\subsection{Deductive Verification}
\label{sub:deductive}

To overcome the limitations of purely automated verification tools, deductive verification frameworks provide an alternative approach. One such framework is Why3~\cite{why3}, which is based on the WhyML~\cite{whyml} specification language and first-order logic. Why3 generates proof obligations that can be discharged by a variety of external theorem provers, allowing verification tasks to be expressed in a structured and modular way.

\subsection{Invariants and Security}
\label{sub:invariants_security}

While the formal verification of CRDTs traditionally focuses on structural convergence, such as ensuring that merge operations form a join-semilattice, the correctness of a distributed application often depends on application-level invariants, including access control policies and other security-related constraints that extend beyond the data type itself.

This challenge has been studied in the context of distributed file systems operating under local-first assumptions~\cite{yanakieva2021access}, where security requirements are commonly expressed in terms of confidentiality, integrity, and availability. Previous work in this area has shown that conventional access control models, such as those inspired by POSIX, give rise to conflicts when replicas operate independently and synchronize asynchronously. These studies establish an impossibility result analogous to the CAP theorem, in this systems that prioritize high availability and tolerate network partitions, it is not possible to simultaneously guarantee confidentiality, integrity, and full availability during conflict resolution, meaning that one of these properties must be relaxed when reconciling concurrent updates.

\section{Verification Tools}
\label{sec:tools}

Verification tools for Replicated Data Types differ in how they balance automation, expressiveness, and proof stability. Initial approaches were based on highly expressive interactive proof assistants, such as Isabelle/HOL~\cite{isabelle}, which offer strong guarantees but require substantial expertise and manual effort. More recent tools adopt SMT-based automation~\cite{porre2023masses} to simplify the verification process at the cost of reduced expressiveness and, in some cases, stability.

\subsection{VeriFx}
\label{sub:verifx}

VeriFx~\cite{porre2023masses} represents a significant step toward making the formal verification of CRDTs accessible. It provides a high-level language, inspired by Scala, in which developers can implement replicated data types using functional abstractions, while automatically verifying properties such as Strong Eventual Consistency through SMT solving with Z3~\cite{z3}.

A key strength of VeriFx lies in its ability to connect specification, verification, and deployment. Designs that are proven correct at the specification level can be translated into executable implementations in languages such as Scala and JavaScript, reducing the gap between formally verified models and practical systems.

Despite these advantages, VeriFx presents several fundamental limitations that restrict its applicability to more complex designs. The tool struggles with recursive algorithms, failing, for instance, to verify the Replicated Growable Array due to the recursive nature of its insertion logic, which exceeds the solver’s ability to reason about unbounded execution paths. In addition, VeriFx exhibits a pronounced sensitivity to syntactic changes, where minor modifications, such as renaming a variable, can cause previously successful proofs to fail or viceversa. Finally, VeriFx does not fully capture the semantic intent of certain conflict resolution policies, as it may prove convergence for a data type even when the verified semantics, such as remove-wins instead of add-wins, do not align with the programmer’s intended behavior.

In addition to these tool-specific limitations, SMT-based verification approaches face significant challenges when reasoning about nested and recursive CRDTs. In such cases, solvers may fail to converge within finite time, forcing developers to resort to manual reasoning techniques, such as structural induction, to establish correctness properties.

\subsection{Why3}
\label{sub:why3}

Why3~\cite{why3} is a platform for deductive program verification based on a first-order specification and programming language called WhyML~\cite{whyml}. Rather than relying on a single automated solver, Why3 serves as an intermediate verification layer that generates proof obligations and dispatches them to multiple external provers, such as Alt-Ergo~\cite{alt-ergo}, Z3~\cite{z3}, and CVC~\cite{cvc}.

Compared to purely SMT-based approaches, Why3 offers greater expressiveness and stability when reasoning about complex data structures. Its deductive verification model is particularly well-suited for nested or recursive definitions, which are difficult to handle using direct SMT encodings. In addition, Why3 enforces a static discipline over mutable state and aliasing, avoiding the need for intricate memory models and resulting in proof obligations that are easier to reason about and maintain.

As a consequence, Why3 provides a more flexible and controlled verification environment than fully automated tools such as VeriFx. While this approach requires more user guidance and interaction, it enables the verification of a broader class of replicated data types~\cite{meirim2019static} and offers a solid basis for addressing limitations related to recursion, modularity, and proof stability.

\endinput
